\chapter{Introducción}
\section{Motivación}

El fútbol es uno de los deportes más populares a nivel mundial, además de constituir un fenómeno social, cultural y económico de gran relevancia, movilizando a millones de aficionados y generando un impacto significativo en la industria del deporte. El mercado global de apuestas deportivas superó los 80.000 millones de dólares en 2022, siendo el fútbol el deporte con mayor volumen de apuestas a nivel mundial. En este contexto, la capacidad de predecir el resultado de los partidos, o al menos de conocer la probabilidad de victoria de cada equipo, representa una herramienta de gran valor tanto para los clubes y sus aficionados como para analistas y operadores de apuestas.

En los últimos años, con el auge del aprendizaje automático, la recolección de datos deportivos y su análisis ha experimentado un crecimiento notable. Sin embargo, la predicción de resultados sigue siendo un problema difícil. Los trabajos de referencia en este campo sitúan la precisión de los mejores modelos en torno al 65--70\% \citep{baboota2019predictive, horvat2020review}, lo que refleja tanto la alta variabilidad inherente al fútbol como las limitaciones de los enfoques que solo utilizan información previa al partido.

Los modelos clásicos de predicción no capturan relaciones no lineales ni incorporan la información que se genera durante el propio encuentro. Las técnicas de \textit{machine learning} abren la posibilidad de construir modelos más complejos, capaces de integrar eventos intrapartido y ajustar las probabilidades de victoria a medida que avanza el juego.

Este trabajo propone aprovechar esa oportunidad desarrollando un sistema que estime, en cada minuto del partido, la probabilidad de victoria local, empate y victoria visitante a partir de datos históricos de eventos reales. Mediante ingeniería de características temporales, modelos de clasificación supervisada y una evaluación por punto de corte temporal, el objetivo es alcanzar o superar el 70\% de precisión que los modelos actuales apenas logran con la información completa del encuentro.

\section{Planteamiento del trabajo}

El planteamiento principal de este trabajo se basa en desarrollar un sistema de predicción de resultados de partidos de fútbol desde un enfoque basado en datos, utilizando técnicas de \textit{machine learning}. Para ello, se usará información histórica de partidos y estadísticas relevantes para entrenar distintos modelos predictivos.

Primeramente, se llevará a cabo la recopilación de datos provenientes de distintas fuentes, que podrán incluir tanto estadísticas finales de partidos como información más detallada basada en eventos y tiempo de juego, en función de su disponibilidad. Posteriormente, se realizará un tratamiento y preprocesado de los datos, con el objetivo de extraer variables relevantes que permitan medir el rendimiento de los equipos y estimar la probabilidad de victoria.

En cuanto a la modelización, se compararán distintos algoritmos de \textit{machine learning}, incluyendo modelos como Random Forest y XGBoost, para determinar cuál ofrece el mejor rendimiento en la predicción de resultados. La evaluación de los modelos se realizará mediante técnicas de validación cruzada y utilizando métricas adecuadas de predicción probabilística, complementadas con métricas tradicionales como la \textit{accuracy} o el \textit{F1-score}. Finalmente, se considerará el uso de redes neuronales como alternativa adicional.

El objetivo principal del trabajo será desarrollar un sistema de predicción que permita estimar de manera precisa la probabilidad de victoria de cada equipo dadas unas circunstancias y métricas de juego definidas. De forma complementaria, se pretende analizar el impacto de las variables y modelos utilizados en cada iteración del proceso de modelización, comparando los resultados y las métricas de rendimiento de los diferentes enfoques.

Finalmente, el alcance del trabajo estará condicionado por la disponibilidad y calidad de los datos empleados, y su efectividad será evaluada considerando la alta aleatoriedad inherente al deporte.

\section{Estructura del trabajo}

El trabajo se organiza en siete capítulos cuyo contenido se resume en la Tabla~\ref{tab:estructura}.

\begin{table}[H]
    \centering
    \small
    \renewcommand{\arraystretch}{1.3}
    \begin{tabularx}{\textwidth}{|c|p{3.8cm}|X|}
        \hline
        \textbf{Cap.} & \textbf{Título}               & \textbf{Contenido principal}                                                                                                                                          \\
        \hline
        1             & Introducción                  & Motivación del trabajo, planteamiento del problema y estructura de la memoria.                                                                                        \\
        \hline
        2             & Contexto y estado del arte    & Revisión de trabajos previos sobre predicción de resultados deportivos con técnicas de \textit{machine learning}, organizados por tipo de enfoque y datos utilizados. \\
        \hline
        3             & Objetivos y metodología       & Definición de los objetivos del proyecto y descripción de los marcos metodológicos adoptados: Scrum para la gestión y CRISP-DM para el ciclo de vida del dato.        \\
        \hline
        4             & Identificación de requisitos  & Descripción de la herramienta web propuesta y especificación de sus requisitos funcionales, no funcionales y de datos.                                                \\
        \hline
        5             & Desarrollo                    & Obtención y análisis exploratorio de los datos, ingeniería de características temporales y entrenamiento de los modelos de clasificación.                             \\
        \hline
        6             & Resultados                    & Comparación de modelos mediante \textit{accuracy}, \textit{F1-score} macro y matrices de confusión, con análisis del rendimiento por punto temporal del partido.      \\
        \hline
        7             & Conclusiones y trabajo futuro & Síntesis de los hallazgos principales, valoración del cumplimiento de objetivos y propuesta de líneas de investigación futura.                                        \\
        \hline
    \end{tabularx}
    \caption{Estructura y contenido de la memoria.}
    \label{tab:estructura}
\end{table}